% GENERATED_BY: oc_core_monograph_translation_draft_pipeline_v1.py
% AI_ASSISTED_TRANSLATION_DRAFT: TRUE
% THEOREM_REVIEW_STATUS: PENDING
% THEOREM_REVIEW_MODE: FORMAL_AI_GATE__CODEX_CLI_CHATGPT
% THEOREM_REVIEW_REVIEWER_ID: 
% THEOREM_REVIEW_AUTHORITY_CLASS: 
% THEOREM_REVIEWED_AT: 
% SOURCE_LANGUAGE: EN
% TARGET_LANGUAGE: RU
% SOURCE_EN_REF: content/18_oc_core_1_3_source_audit.tex
\section{Засвидетельствование источника, судьба теорем и публикационная граница}
\label{sec:oc13-publication-boundary}

Эта глава фиксирует публикационную границу Core~1.3. Ее задача состоит в том,
чтобы определить, какая научная инстанция является канонической, как
обеспечивается засвидетельствование источника, как судьба теорем
распределяется по внешним артефактам и как последующие каналы должны
оставаться согласованными с корпусом источника, SPOT и флагманской
монографией. Она не пересматривает заново формальную доктрину или теоремный
остов; эти задачи относятся к смежным главам и приложениям.

\subsection{Каноническая научная инстанция}

Канонической научной инстанцией выпуска является SPOT-проекция
\occode{releases/oc_core_1_3/editorial/OC_CORE_1_3_SCIENCE_SPOT_latest.json}
(идентификатор поверхности \occode{OC_CORE_1_3_SCIENCE_SPOT}; текущий SHA
зафиксирован в поле \occode{metadata.repo_sha}). Научный корпус источника в
\occode{releases/oc_core_1_3/editorial/science_sources/} является локусом
авторства и засвидетельствования, из которого проецируется этот SPOT.
Исторические свидетельства источника, включая более ранние базовые линии
исправления Core~1.2/Core~1.3, остаются опорами происхождения; они не являются
текущей инстанцией проекции. Нынешняя основная монография -- это флагманская
читательская проекция и сводный справочный том, полученный из этих записей.
Она держит вместе лестницу ориентации, формальную доктрину, теоремный маршрут,
эмпирическую дисциплину, синтез, практическую полезность и слои аудита, но не
подменяет SPOT как конечную инстанцию по объему утверждений, статусу
продвижения или границе доказательства.

Это правило важно потому, что фундаментальные программы со временем накапливают
несколько разных текстовых слоев: исходно-авторские определения и теоремы,
пояснительную прозу, заметки по исправлениям и выпуску, сжатые публичные пакеты
и более поздние операционные проекции. Core~1.3 рассматривает эти слои как
различимые, а не взаимозаменяемые.

\subsection{Засвидетельствование источника}

Core~1.3 в строгом смысле строится с приоритетом источника: флагманские
утверждения, используемые в публичном выпуске, должны оставаться привязанными
к публичным свидетельствам источника, а не только к поздней сводной прозе.
Речь не в архивном фетишизме. Речь в том, чтобы остановить дрейф между пятью
вещами, которые в крупных исследовательских программах часто размываются:

\begin{itemize}
\item теоремоносные утверждения,
\item структурное и интерпретационное объяснение,
\item заметки по исправлениям и выпуску,
\item сжатые публичные пакеты,
\item и более поздние операционные проекции.
\end{itemize}

Когда эта граница размывается, сводное предложение может начать выдавать себя
за саму теорему, а выпускной дашборд -- подменять собой подлинную научную
запись. Core~1.3 предотвращает такой сдвиг, удерживая засвидетельствование
источника привязанным к научному корпусу источника и канонической
SPOT-проекции, а не сводя его к скрытым управляющим файлам.

\subsection{Судьба теорем и дисциплина области действия}

Core~1.3 также требует явного разделения судьбы теорем. Не каждая унаследованная
строка выполняет одну и ту же роль в каждом внешнем артефакте.

\begin{itemize}
\item \occode{THEOREM_NATIVE} строки могут нести положительные
      теоремно-нативные утверждения, когда их цепочка доказательства и
      зависимостей остается точной.
\item \occode{BRIDGE_ONLY} строки могут оставаться научно полезными при явных
      примечаниях об области действия, но их нельзя повышать до
      теоремно-нативного статуса.
\item \occode{FRAME_ONLY} и \occode{EXTENSION} строки принадлежат
      поддерживающим слоям, приложениям, более узким сопутствующим артефактам
      или будущей работе, а не положительному корпусу каждого внешнего текста.
\item \occode{REFUTED} строки исключаются из положительного пути доказательства.
      Их можно упоминать только как отрицательную историю аудита или как запись
      фальсификации, но не как живую поддержку.
\item Материал фронтира или гипотез можно обсуждать только в явно помеченном
      фронтирном языке; его нельзя молча повышать через сводную прозу.
\end{itemize}

Публичное положительное утверждение о замыкании строго ограничено
\occode{CORE_1_3_SCIENCE_ONLY}. Строки \occode{FRAME_ONLY} могут по-прежнему
замыкаться как строки PASS на уровне каркаса внутри поверхности синтеза, но
они не становятся самостоятельными эмпирическими повышениями до
теоремно-нативного статуса. Строки \occode{EXTENSION} и \occode{REFUTED}
остаются вне публичной области PASS по разным причинам: строки
\occode{EXTENSION} -- это фронтирные или поддерживающие линии, ожидающие
более позднего продвижения или отклонения, тогда как строки \occode{REFUTED}
-- это линии отрицательного аудита. Ни один из этих классов не может нести в
этом выпуске положительные утверждения о замыкании.
Во избежание недоразумений строка \occode{FRAME_ONLY} с вердиктом
\occode{PASS} исключается из публичного положительного утверждения о
замыкании, если только более поздний релиз не продвинет ее в
\occode{CORE_1_3_SCIENCE_ONLY} с требуемой теоремно-нативной поддержкой.

Это не отступление. Это условие честной публикации. Монография может оставаться
объемной без риторической инфляции только в том случае, если читатель может
различить, какие строки несут теоремы, какие ограничены, а какие относятся к
поддержке или фронтиру.

\subsection{Связь с журнальным ядром и сопутствующими артефактами}

Поэтому Core~1.3 принимает однонаправленную логику публикации, ориентированной
на читателя:
\[
\begin{gathered}
\text{Научный корпус источника}\\
\text{плюс канонический SPOT}\\
\Downarrow\\
\text{Основная монография}\\
\Downarrow\\
\text{Производные внешние артефакты}
\end{gathered}
\]
Эта диаграмма описывает порядок публикации, а не полный граф происхождения
релизного пакета. Несколько редакторских JSON-поверхностей, сгенерированных
научных глав и матриц выпуска являются прямыми SPOT-проекциями или зеркалами
других управляющих канонических поверхностей. Они не становятся истинными лишь
потому, что основная монография на них ссылается, и они не являются потомками
одной только монографии. Монография является флагманским объектом чтения;
научный корпус источника остается локусом авторства и засвидетельствования, а
SPOT остается канонической инстанцией проекции для релизных поверхностей.

Статья журнального ядра по замыслу уже. Сопроводительные досье, матрицы и
релизные поверхности по замыслу еще уже либо более операциональны. Эти
артефакты могут сжимать, сокращать до подмножества или переорганизовывать
флагманский материал для конкретной задачи, но они не могут расширять научное
утверждение, раздвигать нотацию или выходить за пределы закрепленной здесь
дисциплины судьбы теорем.

Практическое следствие простое. Если какой-либо внешний артефакт,
ориентированный на читателя, утверждает нечто более сильное, чем разрешают
корпус источника и SPOT, то внешний артефакт ошибочен, даже если он короче,
чище или проще в продвижении. Флагманская граница -- это правило публикации,
которое и удерживает этот потолок; она не является дополнительным источником
научного авторитета.

\subsection{Правило последующих проекций}

Та же граница стабилизирует и более позднее использование в рецензировании,
сотрудничестве, продуктовых, деловых и публичных каналах. Разные каналы могут
правомерно цитировать разные части флагманского тома:

\begin{itemize}
\item теоремный маршрут, когда речь идет об ограниченном доказательстве,
\item эмпирический маршрут, когда речь идет о дисциплине продвижения,
\item практический маршрут, когда речь идет об использовании и сравнении,
\item аудиторские приложения, когда речь идет об оспаривании и
      трассируемости.
\end{itemize}

Но они не вправе изобретать вторую научную инстанцию. Засвидетельствование
источника, судьба теорем и публикационная граница существуют для того, чтобы
внешнее использование могло оставаться кратким, не превращаясь в
интерпретативный произвол.

\paragraph{Переход к следующей главе.}
С установленной публикационной границей следующая задача -- внутренняя
согласованность: какие классы утверждений использует флагманский том, почему
иерархия проходит через \(K_{12}\) и как верхние уровни остаются правомерными
без риторической инфляции.

\clearpage

\subsection{External criticism closure route}
\label{sec:oc13-source-audit-external-criticism-closure}

The Stanislav Tsukrov external review package is closed through Appendix~\ref{sec:oc13-external-criticism-closure}. The route records which objections received evidence boundaries, which became proof tasks, and which duplicate revisions were merged. It does not move private review text into the public release and it does not promote a canonical claim.

% SOURCE_EN_REF: content/18_oc_core_1_3_source_audit.tex

\subsection{Source-tradition remediation after external review}
\label{sec:oc13-source-tradition-remediation-v3}

The Stanislav Tsukrov review identified a real scholarly weakness: Core~1.3
already cited physics-adjacent complexity literature, but it did not adequately
name the systems, autopoiesis, identity, complexity-measure, tipping-point, and
formal-geometry traditions in which several OC constructs must be judged. This
release therefore lowers the novelty ceiling and makes the inherited ground
explicit.

The hierarchy and systems-language route is now read against general systems
theory and mathematical systems theory, including von Bertalanffy,
Boulding, Miller, Mesarovic--Takahara, and Klir
\cite{bertalanffy_general_system_theory,boulding_general_systems_theory,miller_living_systems,mesarovic_takahara_general_systems,klir_architecture_systems}.
The rebirth and post-collapse identity route is read against autopoiesis and
identity/persistence debates, including Maturana--Varela, Varela, Parfit, and
Wimsatt \cite{maturana_varela_autopoiesis,varela_biological_autonomy,parfit_reasons_persons,wimsatt_reductionism_levels}.
The complexity and early-warning route is read against effective complexity,
predictive information, integrated information, and critical-transition
literature \cite{gell_mann_effective_complexity,bialek_predictability_complexity,tononi_information_integration,scheffer_critical_transitions,dakos_early_warning}.
Formal differential-geometric machinery is treated as inherited mathematical
tooling rather than as an OC novelty claim; Olver is the current compact source
witness for the jet/Lie-group side of that boundary \cite{olver_lie_groups}.

The repaired claim is narrow: OC contributes a typed cross-level interface,
collapse/residue/rebirth grammar, and release discipline. It does not claim to
replace GST, autopoiesis, complexity theory, identity metaphysics, tipping-point
analysis, or formal differential geometry.
